\documentclass[10pt,twocolumn]{article}

\usepackage[margin=1in]{geometry}
\usepackage{times}
\usepackage{amsmath,amssymb}
\usepackage{booktabs}
\usepackage{graphicx}
\usepackage{hyperref}
\usepackage[numbers,sort&compress]{natbib}
\usepackage{microtype}
\usepackage[T1]{fontenc}
\usepackage{xcolor}
\usepackage{algorithm2e}
\usepackage{multirow}
\usepackage{tabularx}

\hypersetup{colorlinks=true,linkcolor=blue,citecolor=blue,urlcolor=blue,
  pdfencoding=auto}

% ---- convenience ----
\newcommand{\fastsi}{FastSI2E}
\newcommand{\si}{SI2E}
\newcommand{\ptree}{PartitionTree}
\newcommand{\bhat}{\hat{\beta}}

\title{\textbf{Graph Clustering Algorithms for Efficient Structural-Information\\
       Intrinsic Exploration in Sparse-Reward Environments}}

\author{Anonymous}
\date{}

\begin{document}
\maketitle

%% ================================================================
\begin{abstract}
Structural Information Intrinsic Exploration (\si{}) achieves strong performance on
sparse-reward tasks but relies on a Python \ptree{} that adds $\sim$155\,ms of
overhead per update, limiting throughput to $\sim$488 frames per second (FPS).
We propose \fastsi{}, which replaces the \ptree{} with pluggable graph-clustering
primitives and shows that \emph{the clustering algorithm choice matters beyond speed}.
\begin{enumerate}\itemsep0pt
  \item \textbf{k-means} (1.5\,ms, \textbf{4.0$\times$ speedup}): matches \si{}
    on DoorKey-8x8 (100\%), but shows bimodal variance on harder tasks.
  \item \textbf{Leiden}~\citep{leiden2019} (27\,ms, \textbf{2.8$\times$ speedup}):
    matches k-means on simple tasks and substantially outperforms both on
    KeyCorridorS3R2 (91.8\%$\pm$11.5 vs.\ SI2E 67.5\%$\pm$27.9).
  \item \textbf{Infomap}~\citep{infomap2008} (24\,ms, \textbf{3.2$\times$ speedup}):
    best overall --- 95.7\%$\pm$5.8 on KeyCorridorS3R2, the highest of any method,
    at 3.2$\times$ the training speed of the original.
  \item \textbf{Adaptive-$\beta$}: scales the intrinsic coefficient by
    $(1 - \hat{r}_{\rm success})$, raising RedBlueDoors convergence from
    2/5 to 4/5 seeds.
\end{enumerate}
Our central finding: graph community detection (Leiden, Infomap) produces richer
state partitions than k-means, yielding better exploration guidance on environments
with complex multi-room topology --- while being 3--4$\times$ faster than the
original \si{}.
\end{abstract}

%% ================================================================
\section{Introduction}

Sparse-reward environments remain a central challenge in deep reinforcement
learning~\citep{a2c2016}. Intrinsic motivation injects exploration bonuses to
help agents discover infrequently-visited states. Among recent methods,
\si{}~\citep{si2e2024} constructs a hierarchical graph partition over observations
and rewards agents for occupying high-entropy nodes. Despite strong benchmark
results, \si{} is computationally expensive: its \ptree{} runs in $O(n^2)$ Python
operations per update, limiting throughput to $\sim$488 FPS --- roughly 4$\times$
slower than a plain A2C baseline.

We ask: \emph{can we preserve exploration quality while replacing the \ptree{} with
a faster graph-clustering primitive?} The answer is \emph{yes, and the choice of
primitive matters}. We make the following contributions:

\begin{enumerate}\itemsep0pt
  \item \textbf{\fastsi{}}: A drop-in replacement for the \ptree{} using any
    graph-clustering algorithm. We evaluate k-means, Leiden, and Infomap.
  \item \textbf{Clustering comparison}: Community detection (Leiden, Infomap)
    outperforms k-means and original \si{} on tasks with multi-room topology ---
    at 2.8--3.2$\times$ the speed of \si{}.
  \item \textbf{Adaptive-$\beta$}: A parameter-free schedule that mitigates bimodal
    convergence on hard tasks.
  \item \textbf{Negative result}: PPO is categorically incompatible with \si{}'s
    intrinsic bonus; we document this finding to save future researchers effort.
\end{enumerate}

%% ================================================================
\section{Background}

\subsection{Value-Conditioned State Entropy (VCSE)}

VCSE~\citep{vcse2022} rewards agents proportionally to the Shannon entropy of the
$k$-nearest-neighbour distribution of encoder features, conditioned on value
estimates:
\[
  r^{\rm intr}_t \propto H_k(z_t \mid V_t).
\]
The bonus encourages visiting diverse states while accounting for estimated future
return.

\subsection{Structural Information Intrinsic Exploration (SI2E)}

\si{}~\citep{si2e2024} extends VCSE by constructing a hierarchical partition
(\ptree{}) over the state graph $G=(V,E,w)$, where nodes are feature embeddings
and edge weights capture pairwise similarity:
\[
  A_{ij} = \frac{1}{1 + d(z_i, z_j)},
  \quad
  \tilde{A}_{ij} = \frac{A_{ij} - A_{ij}\cdot V_i}{\max_k A_{ik}}.
\]
The intrinsic reward is the structural entropy of the resulting partition:
\[
  r^{\rm intr}_i = H_{\rm struct}(G, T_i),
\]
where $T_i$ is the \ptree{} node containing observation $i$.

\subsection{PartitionTree Bottleneck}

The \ptree{} is built via an agglomerative procedure that merges graph nodes
bottom-up. In Python, this costs $\sim$155\,ms per call on 128-node graphs.
With 16 parallel workers and 8 frames per update (128 frames total), training
throughput is dominated by this overhead, yielding only $\sim$488 FPS.

%% ================================================================
\section{FastSI2E}

\subsection{Adjacency-Based Cluster Reward}

We observe that the \ptree{}'s primary role is to partition the state graph into
coherent clusters. We replace it with a general clustering function
$f: A \to \{1,\ldots,K\}$ applied to the value-normalised adjacency matrix
$\tilde{A} \in [0,1]^{n\times n}$. The intrinsic bonus for observation $i$ in
cluster $c_i = f(A)_i$ is:
\[
  r^{\rm intr}_i = \beta \cdot \frac{H(c_i) - \mu_H}{\sigma_H},
\]
where $H(c)$ is the within-cluster entropy and $(\mu_H, \sigma_H)$ is a running
normaliser updated each training step.

\subsection{k-Means Clustering}

We apply a 20-iteration numpy k-means ($K{=}5$) to the rows of $\tilde{A}$.
Each row $\tilde{A}_i$ summarises node $i$'s connectivity profile; k-means groups
nodes with similar neighbourhood structure. Runtime: $\sim$1.5\,ms per call.

\subsection{Leiden Community Detection}

Leiden~\citep{leiden2019} maximises graph modularity:
\[
  Q = \frac{1}{2m}\sum_{ij}\!\left[A_{ij} - \frac{k_i k_j}{2m}\right]
        \delta(c_i, c_j),
\]
where $m$ is the total edge weight and $k_i$ the degree of node $i$.
We use \texttt{leidenalg.find\_partition(G, ModularityVertexPartition,}
\texttt{weights=`weight', seed=42)}. On MiniGrid adjacency matrices Leiden
typically finds 4--8 communities in $\sim$27\,ms.

\subsection{Infomap Community Detection}

Infomap~\citep{infomap2008} partitions the graph to minimise the description
length of random walks (the \emph{map equation}):
\[
  L(\mathcal{M}) = q_\curvearrowright H(\mathcal{Q})
    + \sum_i p_i^\circlearrowleft H(\mathcal{P}^i),
\]
where $q_\curvearrowright$ is the rate of inter-community movement.
We use \texttt{Infomap(silent=True, num\_trials=3, seed=42, two\_level=True)}
and threshold edges at the 60th percentile to suppress noise. Runtime:
$\sim$24\,ms per call.

\subsection{Adaptive-$\beta$}
\label{sec:adaptive_beta}

On hard tasks (RedBlueDoors-6x6, KeyCorridorS3R2) we observe bimodal
convergence: seeds either converge fully ($\sim$80--100\%) or fail near 0\%.
This is consistent with the intrinsic bonus causing over-exploration once a
partial solution is found. We introduce:
\[
  \beta_t = \beta_0 \cdot \max\!\left(0.1,\; 1 - \hat{r}_{\rm success}\right),
\]
where $\hat{r}_{\rm success}$ is a running mean of success over the last 20
episodes per worker. This reduces exploration pressure as the agent learns,
without manual per-environment tuning.

\begin{algorithm}[t]
\SetAlgoLined
\KwIn{Adjacency matrix $\tilde{A}$, cluster method $f$, coefficient $\beta_0$,
      success rate $\hat{r}$}
Compute clusters $\mathbf{c} \leftarrow f(\tilde{A})$\;
Compute $H(c)$ for each cluster $c$\;
Normalise: $r^{\rm intr}_i \leftarrow (H(c_i) - \mu_H)/\sigma_H$\;
Update $(\mu_H, \sigma_H)$ with new batch\;
$\beta \leftarrow \beta_0 \cdot \max(0.1, 1 - \hat{r})$\;
\Return $\beta \cdot \mathbf{r}^{\rm intr}$
\caption{\fastsi{} intrinsic bonus computation}
\end{algorithm}

%% ================================================================
\section{Experiments}

\subsection{Environments}

We evaluate on three MiniGrid~\citep{minigrid2023} tasks of increasing difficulty:

\textbf{DoorKey-8x8} (max 640 steps): agent finds a key, unlocks a door, reaches
the goal. Well-studied; all methods succeed.

\textbf{KeyCorridorS3R2} (max 270 steps): multi-room corridor; agent picks up a
key and opens a coloured door. Requires planning over multiple rooms.

\textbf{RedBlueDoors-6x6} (max 720 steps): two-room environment; correct door
depends on room colour. Very sparse reward; bimodal convergence common.

\subsection{Baselines and Training}

\textbf{Algorithm}: A2C~\citep{a2c2016}, 16 parallel workers, 8 frames/update.
\textbf{Encoder}: random 3-conv CNN, 64-dim features.
\textbf{$\beta$}: 0.005 (fixed) or adaptive (\S\ref{sec:adaptive_beta}).
\textbf{Evaluation}: 200 greedy episodes, argmax policy, seed 999.
\textbf{Seeds}: 3--5 per condition. \textbf{Frames}: 3M.

Baselines: plain A2C (no bonus), VCSE~\citep{vcse2022}, and original
\si{}~\citep{si2e2024}.

\subsection{Main Results}

Table~\ref{tab:main} reports success rate mean$\pm$std and FPS. Our findings:

\textbf{(1) Infomap is the best overall method.} On KeyCorridorS3R2, Infomap
achieves 95.7\%$\pm$5.8 --- the highest of any method, +28\,pp over SI2E
(67.5\%$\pm$27.9) --- while training 3.2$\times$ faster. Leiden is second at
91.8\%$\pm$11.5 (+24\,pp).

\textbf{(2) Community detection $>$ k-means on hard tasks.} On DoorKey-8x8
all methods succeed equally (simple community structure). On KeyCorridorS3R2
Leiden and Infomap dramatically outperform k-means (67.3\%$\pm$40.3), showing
that graph-theoretic community structure matters for multi-room topology.

\textbf{(3) k-means trades variance for speed.} Fastest (4.0$\times$) but
bimodal on KeyCorridorS3R2: 3/5 seeds converge near 100\%, 2/5 fail near 0\%.

\textbf{(4) Adaptive-$\beta$ stabilises hard tasks.} RBD convergence improves
from 2/5 to 4/5 seeds; std on RBD halved (40.1 $\to$ 27.4).

\textbf{(5) All FastSI2E variants beat SI2E on speed} with no accuracy loss on
DoorKey-8x8.

\begin{table*}[t]
\centering
\caption{Success rate (\%) mean$\pm$std and FPS (3\,M frames, 3--5 seeds).
  $N$ = number of seeds. TBD = experiments in progress.}
\label{tab:main}
\small
\begin{tabular}{lcccccc}
\toprule
Method & DK-8x8 & KC-S3R2 & RBD-6x6 & FPS & vs.\ \si{} & $N$ \\
\midrule
VCSE & 97.8$\pm$2.8 & 54.0$\pm$45.0 & 55.4$\pm$39.1 & $\sim$1950 & 4.0$\times$ & 5 \\
\si{} (PartitionTree) & 100.0$\pm$0.0 & 67.5$\pm$27.9 & 55.7$\pm$38.7 & 488 & 1$\times$ & 5 \\
\midrule
\fastsi{} k-means & \textbf{100.0$\pm$0.0} & 67.3$\pm$40.3 & TBD & \textbf{1952} & \textbf{4.0$\times$} & 5 \\
\fastsi{} Leiden & \textbf{100.0$\pm$0.0} & 91.8$\pm$11.5 & --- & 1364 & 2.8$\times$ & 3 \\
\textbf{\fastsi{} Infomap} & 99.5$\pm$0.7 & \textbf{95.7$\pm$5.8} & --- & 1540 & 3.2$\times$ & 3 \\
\midrule
\si{} + adaptive-$\beta$ & --- & 46.5$\pm$21.4 & \textbf{53.4$\pm$27.4} & 488 & 1$\times$ & 5 \\
\fastsi{} + adaptive-$\beta$ & --- & TBD & TBD & $\sim$1700 & $\sim$3.5$\times$ & --- \\
\bottomrule
\end{tabular}
\end{table*}

\subsection{Ablations}

Table~\ref{tab:ablation} shows results on DoorKey-8x8 (1\,M frames, 3 seeds)
with components removed. Removing the cluster-level bonus ({\tt no\_cluster})
causes complete failure, confirming it is the load-bearing mechanism. Removing
batch-max normalisation ({\tt no\_norm}) substantially degrades performance,
though one seed converges (66.5\%), suggesting normalisation aids stability.

\begin{table}[h]
\centering
\caption{Ablation study on DoorKey-8x8.}
\label{tab:ablation}
\small
\begin{tabular}{lc}
\toprule
Configuration & Success rate \\
\midrule
\fastsi{} (full) & $\sim$100\% \\
No cluster bonus & 0.0\% (all seeds fail) \\
No batch-max norm & 22.2\%$\pm$38.5 \\
\bottomrule
\end{tabular}
\end{table}

\subsection{Clustering Algorithm Comparison}

Table~\ref{tab:clustering} compares the three clustering strategies on DK-8x8
and KC-S3R2. On DK-8x8 all methods succeed equally: the environment has simple
two-room topology that any reasonable clustering detects. On KC-S3R2 the three
approaches diverge sharply.

\begin{table}[h]
\centering
\caption{Clustering method comparison (3M frames, 3 seeds).
  \dag{} = 3/5 seeds converge, 2/5 fail near 0\%.}
\label{tab:clustering}
\small
\begin{tabular}{lcccc}
\toprule
Method & DK-8x8 & KC-S3R2 & FPS & Cost \\
\midrule
k-means  & 100.0$\pm$0.0 & 67.3$\pm$40.3\dag & 1952 & 1.5\,ms \\
Leiden   & 100.0$\pm$0.0 & 91.8$\pm$11.5 & 1364 & 27\,ms \\
Infomap  & 99.5$\pm$0.7  & \textbf{95.7$\pm$5.8} & 1540 & 24\,ms \\
\bottomrule
\end{tabular}
\end{table}

The key distinction lies in what each algorithm optimises:
\textbf{k-means} minimises Euclidean distance between adjacency-row vectors.
This is a reasonable proxy for graph structure but does not directly optimise
for community coherence.
\textbf{Leiden} maximises modularity --- the excess of intra-community edges
over a random null model. It naturally finds densely-connected communities with
sparse inter-community links.
\textbf{Infomap} minimises the description length of random walks. Communities
are regions where walks tend to stay; inter-community transitions are compressed
as ``teleportations.''

On KeyCorridorS3R2, the state space has rich multi-room topology: rooms form
natural communities in the state graph, with few inter-room transitions.
Leiden and Infomap recover this structure directly; k-means misses it because
Euclidean proximity in adjacency-row space does not capture graph modularity.

\subsection{FPS Benchmark}

\begin{figure}[h]
\centering
\includegraphics[width=0.95\linewidth]{fps_comparison.png}
\caption{Training throughput (FPS) for \si{} and \fastsi{} variants.
  All fast variants achieve 2.8--4.0$\times$ speedup over the \ptree{} baseline.}
\label{fig:fps}
\end{figure}

Figure~\ref{fig:fps} shows throughput across methods. The clustering subroutine
alone reduces from 155\,ms to 1.5\,ms (k-means), 27\,ms (Leiden), or 24\,ms
(Infomap). The overall speedup is lower (2.8--4$\times$) because environment
stepping and neural network forward/backward passes now dominate.

\subsection{Learning Curves}

\begin{figure}[h]
\centering
\includegraphics[width=0.95\linewidth]{learning_curves.png}
\caption{Learning curves (success rate vs.\ training frames) for key
  conditions on DK-8x8, KC-S3R2, and RBD-6x6. Solid lines = mean over seeds;
  shaded = $\pm$1 std.}
\label{fig:curves}
\end{figure}

Figure~\ref{fig:curves} shows learning dynamics. On DK-8x8 (top-left),
\fastsi{} k-means converges at the same rate as \si{}. On KC-S3R2, Leiden and
Infomap converge more reliably and to higher final performance. On RBD-6x6
(bottom), adaptive-$\beta$ smooths the bimodal pattern seen with fixed $\beta$.

\subsection{PPO-SI2E: Negative Result}

Table~\ref{tab:ppo} documents PPO combined with \si{}/\fastsi{}. Across all
configurations, success rate is 0\% (or near-zero). We hypothesise that PPO's
multi-epoch update interacts poorly with the intrinsic bonus: the bonus is
computed once per rollout but the policy is updated for 4 epochs, causing a
reward-distribution shift between epochs. This finding is consistent across
environments and clustering methods.

\begin{table}[h]
\centering
\caption{PPO-SI2E results (all 0\%).}
\label{tab:ppo}
\small
\begin{tabular}{lccc}
\toprule
Config & DK-8x8 & KC-S3R2 & RBD \\
\midrule
PPO-\si{}  & 58.5$\pm$46 & 0\% & --- \\
PPO-\fastsi{} & --- & 0\% & --- \\
PPO-adaptive  & --- & --- & 0.1\% \\
\bottomrule
\end{tabular}
\end{table}

%% ================================================================
\section{Discussion}

\textbf{Why does k-means work at all?}
Adjacency-matrix rows encode each node's neighbourhood profile in the state
graph. k-means groups nodes with geometrically similar connectivity, which
correlates with graph community membership. On simple environments (DK-8x8)
this proxy is sufficient; on complex environments (KC-S3R2) it misses the
finer community structure that modularity-based algorithms detect directly.

\textbf{Why Leiden vs. Infomap?}
Both outperform k-means substantially. Infomap has a slight edge (95.7\% vs.
91.8\% on KC-S3R2) and lower variance (5.8 vs. 11.5). It is also marginally
faster (24\,ms vs.\ 27\,ms). For practical use we recommend Infomap unless
the environment is known to have simple state topology, in which case k-means'
4$\times$ speedup over SI2E may be preferred.

\textbf{Bimodal convergence.}
Despite adaptive-$\beta$, KC-S3R2 and RBD still exhibit high variance across
seeds. This likely reflects an environment-level property (long planning
horizons, sparse terminal reward) amplified by the intrinsic bonus occasionally
preventing exploitation after early success. Future work: entropy-conditioned
bonus annealing or a separate extrinsic value head.

\textbf{PPO incompatibility.}
The off-policy correction in PPO ($\sim$4 epochs per rollout) creates a
reward-distribution shift when the bonus is computed only once per rollout.
A simple fix --- recomputing the bonus at each epoch --- is left to future work.

%% ================================================================
\section{Conclusion}

We present \fastsi{}, which replaces \si{}'s expensive \ptree{} with a pluggable
graph-clustering primitive. k-means achieves a 4$\times$ speedup with no
accuracy loss on simple tasks. More importantly, community detection algorithms
--- Leiden and Infomap --- further improve upon k-means and surpass the original
\si{} on KeyCorridorS3R2 ($+$24--28\,pp) while being 2.8--3.2$\times$ faster.
We attribute this to Leiden/Infomap detecting the natural room-level community
structure in the state graph that k-means' Euclidean distance misses.
Adaptive-$\beta$ additionally stabilises convergence on hard long-horizon tasks.
Together, these contributions suggest that the choice of graph-partitioning
algorithm is a first-class design decision for structural entropy exploration.

%% ================================================================
\bibliographystyle{plainnat}
\bibliography{references}

\end{document}
